\section{Introduction}

\label{sec:1}

Homogeneous geometries provide a finite set of variables in which the dependence of a gravitational action on size, shape and torsion can be studied explicitly. In Einstein--Cartan (EC) theory the coframe and a metric-compatible connection allow these ingredients to be varied separately \cite{Hehl1976}. The complete Nieh--Yan (NY) term fixes how the torsion and curvature contributions enter the same differential-form identity \cite{NiehYan1982}. On $S^3\times S^1$, spatial rotations organize the geometric variables into scalar, vector and traceless-tensor representations.

Paper01 classified radius extrema on fixed homogeneous families, while Paper02 established curvature identities and shape bounds including an irreducible EC-Weyl correction \cite{paper1,paper2}. Here we study the EC member of this family, with the curvature-squared coefficient set to zero, and develop its spin-0, spin-1 and spin-2 geometric sectors. The term \emph{spin} refers to the spatial $SO(3)$ representation of a deformation. We use it to organize variables and interactions; physical propagation requires a separate dynamical and constraint analysis.

The three sectors organize both a static potential and an exact circle-dependent action. Spin-0 gives a conditional radius minimum, radius--torsion couplings and a radial derivative term. Spin-1 has a positive static twist cost and, at fixed scalar controls and round shape, a finite-amplitude derivative term proportional to $|\omega\times\omega'|^2$. Spin-2 gives a continuous bifurcation of constrained static minima, a square-root amplitude law and a second-derivative jump; minimizing the static potential over constant volume-preserving shapes and twist at fixed scalars selects the round, zero-twist geometry. At zero twist and fixed radius, the circle variation supplies a nonnegative shape derivative term.

The joint circle action connects these results through derivative cross terms and the spatial shift. At the round, zero-twist reference, the derivative quadratic form separates the sectors. Nonlinear cross terms couple them, and an explicit periodic shape--twist family makes the joint Euclidean action unbounded below even at fixed scalar controls. Static shape minima and minima of the circle functional are consequently different variational questions. Section 2 gives the general action, Sections 3--5 develop the three sectors, and Section 6 derives the interactions and coupled family. Section 7 relates these results to Papers01/02.